\section{Project Objectives Decomposition}
\label{sec:project-objectives-decomposition}

The project objectives defined in Section \ref{subsec:aims-and-objectives} can be decomposed into several smaller tasks which allow for a more detailed plan to be created.
These areoutlined in Table \ref{tab:project-tasks} below.
This tasks span three main contributory areas:
(1) creation of the gait-based identification system (\ac{grs}),
(2) creation of the access control system demonstration (\ac{acc}),
(3) testing and evaluation of the final product.

\begin{scriptsize}
    \begin{longtable}[c]{|p{0.03\textwidth}|p{0.7\textwidth}|p{0.17\textwidth}|}
        \caption{Project Task List}
        \label{tab:project-tasks}                                                                                                                                                                                                                                                                                                                                                                                                                                      \\
        \hline
        \textbf{\#} & \textbf{Task}                                                                                                                                                                                                                                                                                                                                                                                                            & \textbf{Dependencies} \\ \hline
        \endfirsthead
        \caption*{Project Task List (Continued)}                                                                                                                                                                                                                                                                                                                                                                                                                       \\
        \hline
        \textbf{\#} & \textbf{Task}                                                                                                                                                                                                                                                                                                                                                                                                            & \textbf{Dependencies} \\ \hline
        \endhead
        \hline
        \endfoot
        1           & \textbf{Identify a suitable gait-data extraction method and \ac{gr} approach to be implemented.}                                                           \par From the approaches discussed in Section \ref{sec:gait-recognition-for-physical-access-control}, a suitable option should be chosen that will allow for completion of a \ac{poc} within the timeframe of this project, with hardware that is accessible. &                       \\ \hline
        2           & \textbf{Acquire the required hardware for the chosen \ac{gr} approach.}                                                        \par A \ac{cv}-based approach will be used, but the chosen method may influence the type of camera required and the computing power of the device running the application(s).                                                                                                             & 1                     \\ \hline
        3           & \textbf{Identify the required software resources to implement the chosen \ac{gr} approach.}                           \par A programming language and required libraries should be chosen that will allow for implementation of the this systems functionality, this will likely be Python.                                                                                                                              & 1                     \\
        4           & \textbf{Produce code to collect movement data using \ac{cv} and abstract it into a suitable representation.}                          \par Utilising the acquired hardware and chosen programming language, functionality to capture movement data from a camera feed and abstract it into a suitable gait cycle representation should be implemented.                                                                   & 1, 2, 3               \\ \hline
        5           & \textbf{Produce \ac{grs} code to compare input gait cycle data (Task \#4) to historic data and predict identity.}                           \par Code should be created that compares fresh input data to historic data in order to classify it and attempt to identify the associated individual.                                                                                                                       & 4                     \\ \hline
        6           & \textbf{Migrate gait cycle collection code (Task \#4) to the \ac{acc} codebase.}                                                \par The \ac{acc} will ultimately be responsible for collecting gait cycle data in real-time, therefore this functionality should be moved here.                                                                                                                                         & 4                     \\ \hline
        7           & \textbf{Produce code to allow the \ac{acc} and \ac{grs} to communicate.}                                         \par To allow the \ac{acc} to send gait data to and receive responses from the \ac{grs}, network communication functionality must be implemented.                                                                                                                                                       &                       \\ \hline
        8           & \textbf{Create structure of the \ac{grs}' access control rule database.}                                                             \par As shown in Figure \ref{fig:system-topology}, the \ac{grs} will use a database that stores access control rules to determine which identities have access.                                                                                                                     &                       \\ \hline
        9           & \textbf{Produce code to check access control rules and permit or deny entry based on the identity returned by the \ac{grs}.}                             \par The \ac{grs} must be able to query the database created in Task \#8 and retrieve the rule associated with a given identity. This will be returned to the \ac{acc} with identificaiton results.                                                             & 8                     \\ \hline
        10          & \textbf{Create a \ac{acc} \ac{ui} to allow the registration of gait-based identities.}                         \par To boost usability enable gait-based identity registration a suitable user interface (e.g., graphical) should be created for the \ac{acc}.                                                                                                                                                           &                       \\ \hline
        11          & \textbf{Add a form of binary output to the \ac{acc} interface that demonstrates access being permitted or denied based on identification.}                              \par Either physical or digital output should be implemented to demonstrate the provision or denial of access based on identification results returned by the \ac{grs}.                                                                          & 10                    \\ \hline
        12          & \textbf{Create a \ac{grs} \ac{ui} for managing registered identities and associated access control rules.} \par The \ac{ui} should allow administrators to view and delete registered identities, as well as manage their associated access control rules. This will facilitate easier management of the system and ensure that access can be dynamically adjusted as needed.                                            &                       \\ \hline
        13          & \textbf{Create or acquire a historic dataset relevant to the chosen \ac{gr} approach.}                                               \par Depending on the chosen approach (e.g., silhouette or model-based, and \ac{ml} or model-free), a dataset must be acquired or created to enable reliable testing.                                                                                                               & 1                     \\
        14          & \textbf{Create historic gait data store (and/or train a predictive \ac{ml} model) using Task \#13 data.}                           \par For a model-free approach, the acquired dataset should be formatted (e.g., in a folder) suitably, or if \ac{ml} is being used, a model must be trained using the data that will allow for subsequent classification.                                                             & 13                    \\ \hline
        15          & \textbf{Complete testing of the product of Tasks \#1-13 using Task\#14 data in respect to Objective \#3.}                                                         \par The final product should be adequately tested to ensure it is capable of identifying individuals based on gait cycle both efficiently and accurately.                                                                                             & 1-14                  \\ \hline
        16          & \textbf{Perform demonstrations and collect feedback and opinions on the resulting \ac{gr}-based access control system.}                                                           \par Receiving feedback where possible is important, as one of the goals is to improve user experience with \ac{pacs} authentication.                                                                                                  & 15                    \\ \hline
        17          & \textbf{Evaluate the final product, address problems, and identify areas for improvement.}                                                 \par Based on self-assessment and feedback received from others, assess how the project has met the initially defined goals, and where improvements could be made if developed further.                                                                                       & 16                    \\ \hline
        18          & \textbf{Discuss the feasibility of implementation in a real-world setting and what must change to make that happen.}                          \par After producing a Proof-of-Concept Gait-based authentication system, it should be evaluated to determine how it could be implemented in a real-world setting, associated risks, and changes that may need to be made to improve robustness or change functionality.   & 17                    \\
    \end{longtable}
\end{scriptsize}

An agile development approach will be taken during implementation of this \ac{poc}.
This will enable the separation of concerns, and for tasks to be prioritised depending on what is needed to progress.
It will also allow for adjustments to be made to the project approach if needed, and for tasks to be worked on in parallel, such as elements of the \ac{acc} and \ac{grs} that are dependent on one another.
